\documentclass[10pt,a4paper,onecolumn]{article}

\usepackage{amsmath,amssymb}
\usepackage{graphicx}
\usepackage[margin=2cm]{geometry}
\usepackage{hyperref}
\usepackage{cite}
\usepackage{enumitem}
\usepackage{booktabs}
\usepackage{indentfirst}
\usepackage[bahasa]{babel}

\hypersetup{colorlinks=true,linkcolor=black,urlcolor=blue,citecolor=blue}
\setlength{\parskip}{0.15ex}
\setlength{\parindent}{1.25cm}

\title{\textbf{Segmentasi Deforestasi Multi-Temporal \\
Menggunakan Sentinel-2 Level-2A dan Label Hansen GFW} \\
\small Pipeline Ekspor--Tile--Label--Training untuk Hutan Lindung Papua Selatan}

\author{Tim Riset --- Deforest.id}
\date{}

\begin{document}
\maketitle
\thispagestyle{empty}

\paragraph{Abstrak.}
Kami menyajikan pipeline otomatis untuk segmentasi deforestasi pada
resolusi 10\,m di Hutan Lindung Boven Digoel dan Merauke, Papua Selatan.
Pipeline mencakup: (1) ekspor komposit Sentinel-2 Level-2A tahun 2020
dan 2023 dari Google Earth Engine, (2) pemotongan menjadi chip
$64\times64$ piksel dengan filter tutupan awan dan NDVI, (3) pelabelan
menggunakan Hansen Global Forest Change v1.11, dan (4) pelatihan model
U-Net. Tiga arsitektur dievaluasi: SimpleUNet satu tanggal, ResNet34
U-Net multi-temporal, dan FC-Siam-diff. Model Siamese mencapai IoU uji
global 0,137 (F1 0,240), melampaui multi-temporal (0,119) dan satu
tanggal (0,122). Ketimpangan kelas ekstrem (0,38\,\% piksel positif) dan
ketidaksesuaian resolusi label Landsat 30\,m--Sentinel-2 10\,m membatasi
performa absolut. Pipeline menunjukkan kelayakan ujung-ke-ujung dari
akuisisi hingga pelatihan model segmentasi.

\section{Pendahuluan}
Indonesia memiliki sekitar 94~Mha hutan tropis dan kehilangan ratusan
ribu hektare hutan setiap tahun \cite{hansen2013high}. Hutan lindung di
Papua Selatan, khususnya Kabupaten Boven Digoel dan Merauke, merupakan
kawasan yang rentan terhadap pembukaan lahan skala kecil yang sulit
terdeteksi oleh sistem resolusi kasar. Global Forest Watch (GFW)
menyediakan data kehilangan tutupan pohon tahunan pada resolusi 30~m,
sementara Sentinel-2 pada resolusi 10~m memberikan dasar praktis untuk
deteksi yang lebih rapat \cite{drusch2012sentinel2}.

Studi ini membangun pipeline ujung-ke-ujung yang mencakup: (1) ekspor
komposit Sentinel-2 Level-2A (\texttt{COPERNICUS/S2\_SR\_HARMONIZED})
dari Google Earth Engine \cite{gorelick2017gee} untuk periode baseline
2020 dan deforestasi 2023, (2) pemotongan menjadi chip $64\times64$ pada
resolusi 10~m dengan filter tutupan awan dan NDVI, (3) pelabelan
menggunakan Hansen Global Forest Change (GFC) v1.11 \cite{hansen2013high},
dan (4) pelatihan tiga arsitektur segmentasi. Kami menguji apakah label
loss-year Hansen GFW dapat digunakan untuk mengawasi segmentasi
Sentinel-2 dan apakah pemodelan multi-temporal meningkatkan kinerja
dibandingkan segmentasi satu tanggal pasca-kejadian.

\section{Kajian Terkait dan Posisi Riset}
Produk kehilangan hutan tahunan Hansen GFC \cite{hansen2013high} berbasis
Landsat (30~m) tetap menjadi acuan global, tetapi granularitas tahunan
dan resolusi spasialnya membatasi deteksi gangguan kecil. Sentinel-2
menawarkan resolusi 10~m dengan revisit 5 hari, cocok untuk pemantauan
lokal \cite{drusch2012sentinel2}. U-Net \cite{ronneberger2015unet}
banyak digunakan untuk segmentasi penginderaan jauh, termasuk deteksi
deforestasi \cite{wagner2019change,john2022attention}. Arsitektur
change detection seperti FC-Siam-diff \cite{daudt2018fcsiam}
membandingkan fitur dua tanggal secara eksplisit dan relevan untuk
pemantauan deforestasi multi-temporal.

Riset ini berfokus pada implementasi pipeline operasional yang nyata,
bukan sekadar tolok ukur akurasi. Kontribusi utamanya adalah mengukur
kesenjangan antara label Hansen 30~m dan citra Sentinel-2 10~m, serta
mengidentifikasi hambatan pipeline: resolusi label, ketimpangan kelas,
overfitting, dan artefak komposit akibat tutupan awan.

\section{Metodologi}
\subsection{Data dan Label}
Wilayah kajian adalah Hutan Lindung di Kabupaten Boven Digoel dan
Merauke, Provinsi Papua Selatan (longitude $\sim$138,1--139,9\,E,
latitude $\sim$5,9--4,5\,S). Tujuh area sampel (\texttt{hl\_sample\_1}
hingga \texttt{hl\_sample\_7}) didefinisikan dalam berkas GeoJSON
berdasarkan data kawasan hutan lindung resmi.

Kami mengekspor komposit Sentinel-2 Level-2A
(\texttt{COPERNICUS/S2\_SR\_HARMONIZED}) dari Google Earth Engine
\cite{gorelick2017gee} untuk dua periode: baseline 2020
(2020-01-01--2020-12-31) dan deforestasi 2023 (2023-01-01--2023-12-31).
Setiap komposit dibangun dengan filter \texttt{CLOUDY\_PIXEL\_PERCENTAGE}
$<80$, masking awan dan cirrus via QA60 (bit 10 dan 11), dan median
per-piksel. Lima kanal diekspor: B2 (Biru), B3 (Hijau), B4 (Merah), B8
(NIR), dan \texttt{CLEAR\_COUNT} --- jumlah observasi jernih per piksel.
Resolusi spasial 10~m, CRS EPSG:4326.

Komposit GeoTIFF dipotong menjadi chip $64\times64$ piksel dengan
stride 64, setara 0,64~ha per chip. Filter diterapkan: chip dengan
fraksi awan $>0{,}3$ (berdasarkan \texttt{CLEAR\_COUNT} $<3$) atau
NDVI rata-rata $<0{,}0$ dibuang. Setiap chip menyimpan larik
\texttt{rgb} (3 kanal ternormalisasi), \texttt{nir}, \texttt{ndvi}
(hitungan dari B4 dan B8), dan \texttt{cloud} (masker biner).

Label lemah diturunkan dari Hansen Global Forest Change v1.11
\cite{hansen2013high}. Piksel dengan \texttt{loss} $>0$ pada tahun
2023 diberi label deforestasi (1); piksel lain diberi label latar
belakang (0). Masker GFW diunduh dari \texttt{storage.googleapis.com}
dan diresampel ke resolusi 10~m dengan interpolasi nearest-neighbor.
Piksel awan pada chip diberi nilai 255 dan diabaikan saat kalkulasi
loss (ignore index).

Dataset berisi 40.484 chip dari 7 scene: 23.135 latih, 5.822 validasi,
11.527 uji. Hanya 0,38\,\% dari seluruh piksel yang positif dan 74\,\%
chip tidak mengandung deforestasi, mencerminkan ketimpangan kelas
ekstrem yang khas dalam pemantauan operasional.

\subsection{Model dan Pelatihan}
Tiga arsitektur dievaluasi. Pertama, \textbf{SimpleUNet} \cite{ronneberger2015unet}
hanya menggunakan citra pasca-kejadian 2023 dengan lima kanal input
(rgb+nir+ndvi). Kedua, \textbf{ResNet34 U-Net} \cite{he2016deep}
menggabungkan input pra-kejadian 2020 dan pasca-kejadian 2023 menjadi
10 kanal, menggunakan encoder ResNet34 yang diinisialisasi dengan bobot
ImageNet (tiga kanal pertama dari bobot RGB, tujuh kanal tambahan dari
rata-rata bobot RGB). Ketiga, \textbf{FC-Siam-diff} \cite{daudt2018fcsiam}
menggunakan encoder berbobot sama untuk dua tanggal dan menggabungkan
fitur pasca-kejadian, fitur pra-kejadian, serta selisih fiturnya pada
skip connection decoder.

Semua model menggunakan Dice loss \cite{milletari2016vnet} dengan
\texttt{ignore\_index}=255 (piksel awan), AdamW \cite{loshchilov2019decoupled}
(laju belajar $10^{-4}$, weight decay $5\times10^{-4}$), cosine annealing
scheduler, dan \texttt{WeightedRandomSampler} dengan faktor 3--10$\times$
untuk chip positif guna mengatasi ketimpangan kelas. SimpleUNet dilatih
dengan batch 32--64, sedangkan model multi-temporal menggunakan batch 64
dengan gradient accumulation 2 langkah (batch efektif 128). Augmentasi
terbatas pada flip horizontal/vertikal acak; augmentasi berat seperti
rotasi dan jitter spektral justru menurunkan hasil karena model bergantung
pada hubungan kanal yang konsisten. Mixed precision (AMP) dan pemotongan
gradien 1,0 digunakan pada semua pelatihan.

Evaluasi melaporkan IoU, precision, recall, F1, dan akurasi karena
akurasi saja dapat menyesatkan ketika latar belakang mendominasi
\cite{powers2011evaluation}. Pada inferensi, test-time augmentation
horizontal dan vertikal merata-ratakan prediksi untuk mengurangi
sensitivitas orientasi chip.

\section{Hasil dan Diskusi}
\begin{table}[ht]
\centering
\caption{Metrik global pada test set, dijumlahkan pada semua chip.}
\label{tab:results}
\small
\begin{tabular}{lccc}
\toprule
Metrik & SimpleUNet & ResNet34 UNet & Siamese \\
\midrule
IoU       & 0,122 & 0,119 & \textbf{0,137} \\
Precision & \textbf{0,198} & 0,154 & 0,196 \\
Recall    & 0,239 & \textbf{0,344} & 0,311 \\
F1        & 0,217 & 0,212 & \textbf{0,240} \\
Accuracy  & \textbf{0,984} & 0,977 & 0,982 \\
\bottomrule
\end{tabular}
\end{table}

Model Siamese mencapai IoU dan F1 terbaik, mengonfirmasi bahwa perbedaan
fitur eksplisit lebih berguna daripada penumpukan kanal. ResNet34 UNet
memiliki recall tertinggi ({0,344}) tetapi precision lebih rendah, yang
menunjukkan kecenderungan positif palsu. SimpleUNet satu tanggal justru
paling presisi ({0,198}) namun melewatkan lebih banyak kejadian. Rata-rata
IoU per sampel jauh lebih rendah ({0,012--0,015}) karena mayoritas chip
tidak mengandung piksel positif; metrik global lebih representatif untuk
kejadian langka.

Kurva validasi mencapai puncak pada epoch 4--6 lalu menurun sementara
loss pelatihan terus turun, menandakan overfitting akibat contoh positif
terbatas dan konteks $64\times64$ yang sempit. Analisis kesalahan
menunjukkan empat mode kegagalan utama: positif palsu dari perubahan
pertanian/tanah terbuka/air, batas deforestasi yang terlalu halus,
artefak bayangan awan residual, dan noise label spasial/temporal akibat
ketidakcocokan resolusi GFW 30~m--Sentinel-2 10~m.

\section{Pipeline dan Alur Kerja}
Pipeline diimplementasikan dalam modul \texttt{services/gee-export} dan
\texttt{scripts/} dengan CLI terpadu:

\begin{enumerate}[nosep]
    \item \textbf{Ekspor} --- \texttt{python -m src.run export --geojson aoi.geojson --prefix hl\_sample} \\
    Mengekspor komposit Sentinel-2 L2A (lima kanal: B2, B3, B4, B8, CLEAR\_COUNT) dalam format GeoTIFF ke Google Drive atau Cloud Storage.
    \item \textbf{Tiling} --- \texttt{python -m src.run tile} \\
    Memotong GeoTIFF menjadi chip $64\times64$ dengan stride 64. Filter:
    fraksi awan $<0{,}3$ (CLEAR\_COUNT $\geq3$) dan NDVI rata-rata $>0{,}0$.
    Setiap chip disimpan sebagai \texttt{.npz} dengan larik rgb, nir, ndvi, dan cloud.
    \item \textbf{Labeling} --- \texttt{python -m src.run label-gfw} \\
    Mengunduh tile Hansen GFC-2023-v1.11 dari Google Cloud Storage dan
    meresampel masker loss (30~m $\to$ 10~m) pada setiap chip.
    \item \textbf{Split} --- \texttt{python -m src.run split} \\
    Membagi dataset menjadi train/val/test (70/20/10) dan menyimpan dalam
    format \texttt{.npy} siap pakai untuk PyTorch.
    \item \textbf{Training} --- \texttt{python scripts/train\_unet.py} \\
    Melatih SimpleUNet (5 kanal) atau varian multi-temporal dengan Dice
    loss, AdamW, WeightedRandomSampler, dan cosine annealing scheduler.
\end{enumerate}

Pipeline berjalan di atas GPU dengan PyTorch AMP. Pelatihan 30 epoch
membutuhkan kurang dari 30 menit; inferensi penuh pada 11.527 chip uji
dengan test-time augmentation membutuhkan sekitar 12 menit. Seluruh
alur kerja dapat direproduksi pada perangkat keras konsumen.

Keterbatasan utama adalah kualitas label Hansen GFC yang berasal dari
Landsat 30~m, sehingga setiap piksel label mencakup sembilan piksel
Sentinel-2. Pembukaan kecil yang terlihat pada 10~m dapat hilang dari
label. Ketimpangan kelas ekstrem (0,38\,\% positif), konteks chip
$64\times64$ yang terbatas, artefak komposit, dan pergeseran distribusi
antar-scene juga membatasi generalisasi. Keterbatasan ini menunjukkan
bahwa peningkatan kualitas label lebih berdampak daripada perubahan
arsitektur model.

\section{Kesimpulan}
Pipeline ujung-ke-ujung ini menunjukkan bahwa segmentasi deforestasi
10~m menggunakan Sentinel-2 Level-2A dengan label lemah Hansen GFC v1.11
layak diimplementasikan untuk hutan lindung di Papua Selatan. FC-Siam-diff
mencapai hasil terbaik (IoU 0,137, F1 0,240), mengonfirmasi bahwa
representasi perubahan eksplisit mengungguli stacking naif dan inferensi
satu tanggal. Pekerjaan lanjutan perlu berfokus pada label berkualitas
lebih tinggi (manual atau citra resolusi tinggi), chip lebih besar dengan
mekanisme attention, dan pascapemrosesan komponen terhubung untuk
peringatan operasional. Seluruh pipeline --- dari ekspor GEE hingga
pelatihan model --- bersifat terbuka dan dapat direproduksi.

\clearpage
\small
\renewcommand{\refname}{Daftar Pustaka}
\bibliographystyle{ieeetr}
\bibliography{icp-refs}

\clearpage
\appendix
\section{Arsitektur SimpleUNet}
\label{app:arch}
Encoder--decoder dengan $5\to32\to64\to128\to256$ kanal pada encoder,
bottleneck 512, dan skip connection. Input: 5 kanal (rgb ternormalisasi
0--255, nir/10000, ndvi). Output: 2 kanal (background, deforestasi).
Loss: Dice dengan ignore\_index=255.

\end{document}